\chapter{Transition Systems and Temporal Properties}
\label{ch:transition-temporal-properties}

A code cache can contain correct machine instructions and still execute freed memory.
The missing invariant concerns the order of publication, invalidation, and reclamation.
Likewise, a deoptimization routine can reconstruct each value correctly yet resume execution before all required values exist.
These failures concern sequences of states, including choices made by other runtime components.
A transition model makes those choices explicit; a temporal property states which sequences are permitted.
The resulting proof concerns the model, so deciding what a state remembers is as important as choosing a checking algorithm.

\section{A state records enough to predict the next step}
\label{sec:ch45-state-model}

Consider one compiled code-cache entry and one possible user of it.
We abstract away instruction bytes and retain seven protocol states.
State $c_0$ is empty; $c_1$ is emitting code; $c_2$ has code and reconstruction metadata ready; $c_3$ is published and pinned by a user; $c_4$ is unpublished but still pinned; $c_5$ is reclaimed; and $c_6$ is reclaimed while a user still holds a pin.
A pin represents a claim that prevents reclamation.
We model one compilation and one use, with no reuse of the entry after reclamation.
The model therefore says nothing yet about address reuse, generation counters, or two simultaneous users.

Figure~\ref{fig:ch45-cache} gives the transitions.
The self-loop at $c_1$ represents an emission step that has not finished; the loops at $c_3$ and $c_4$ represent a user that has not released its pin.
In this deliberately small model, publication and acquisition of the first pin form one atomic step.
A model of the actual publication protocol would separate these events if another thread could intervene.
The dashed edge represents a proposed, faulty reclamation rule, and is absent from the baseline model.

\begin{figure}[tbp]
\centering
\begin{tikzpicture}[node distance=10mm and 13mm]
\node[basic block] (c0) {$c_0$: empty};
\node[basic block,right=of c0] (c1) {$c_1$: emit};
\node[basic block,right=of c1] (c2) {$c_2$: ready};
\node[basic block,right=of c2] (c3) {$c_3$: live};
\node[basic block,below=16mm of c3] (c4) {$c_4$: drain};
\node[basic block,left=of c4] (c5) {$c_5$: freed};
\node[basic block,left=of c5] (c6) {$c_6$: error};
\draw[cfg edge] ([xshift=-8mm]c0.west) -- (c0);
\draw[cfg edge] (c0) -- (c1);
\draw[cfg edge] (c1) -- (c2);
\draw[cfg edge] (c2) -- (c3);
\draw[cfg edge] (c3) -- (c4);
\draw[cfg edge] (c4) -- node[above,font=\scriptsize] {release} (c5);
\draw[back edge] (c4.south) to[bend left=35] node[below,font=\scriptsize] {faulty reclaim} (c6.south);
\draw[cfg edge] (c1) to[loop above] (c1);
\draw[cfg edge] (c3) to[loop above] (c3);
\draw[cfg edge] (c4) to[loop right] (c4);
\draw[cfg edge] (c5) to[loop above] (c5);
\draw[cfg edge] (c6) to[loop left] (c6);
\end{tikzpicture}

\caption{The baseline cache protocol waits for the last pin; the dashed candidate edge reclaims too early.}
\label{fig:ch45-cache}
\end{figure}

Formally, a finite \term{Kripke structure} is $K=(S,I,R,AP,L)$.
The finite, nonempty set $S$ contains states; the nonempty subset $I\subseteq S$ contains allowed initial states; $R\subseteq S\times S$ is the transition relation; $AP$ is a finite set of atomic propositions; and $L:S\to 2^{AP}$ gives the propositions true at each state.
We require $R$ to be total: every state has at least one successor.
Section~\ref{sec:ch45-deadlocks} explains how to satisfy that requirement without silently treating a stuck execution as successful termination.
An edge represents one modeled step, not a unit of physical time.

An \term{atomic proposition} is an observable Boolean fact, such as $\mathit{pinned}$ or $\mathit{error}$.
It is not an arbitrary name attached to an edge.
For this example $I=\{c_0\}$, and Table~\ref{tab:ch45-cache-labels} computes the relevant labels.
The proposition $\mathit{freed}$ distinguishes completed reclamation from the initial absence of storage.
In the error state, $\mathit{pinned}$ remains true because the outstanding user's claim is exactly what makes reclamation invalid.

\begin{table}[tbp]
\centering
\caption{State labels expose the difference between absence of code and premature reclamation.}
\label{tab:ch45-cache-labels}
\begin{tabular}{lrrrrrr}
\toprule
State & resident & metadata & published & pinned & freed & error \\
\midrule
$c_0$ & 0 & 0 & 0 & 0 & 0 & 0 \\
$c_1$ & 1 & 0 & 0 & 0 & 0 & 0 \\
$c_2$ & 1 & 1 & 0 & 0 & 0 & 0 \\
$c_3$ & 1 & 1 & 1 & 1 & 0 & 0 \\
$c_4$ & 1 & 1 & 0 & 1 & 0 & 0 \\
$c_5$ & 0 & 0 & 0 & 0 & 1 & 0 \\
$c_6$ & 0 & 0 & 0 & 1 & 1 & 1 \\
\bottomrule
\end{tabular}
\end{table}

A CFG vertex is usually a program point shared by many executions.
A transition-system state can instead contain $(\mathit{pc},\mathit{store},\mathit{phase},\mathit{pending})$.
Two executions at the same instruction with different pin counts are distinct states if their legal next steps differ.
Conversely, several machine instructions may be summarized by one edge if the abstraction justifies that summary.
Calling a CFG a transition system does not supply the missing store or prove that an infeasible CFG path is executable.

Finite source code does not imply a finite state space.
Unbounded heaps, mathematical integers, recursion, and unbounded thread creation all require bounds or abstraction.
Fixed-width values and a fixed memory bound give a finite model, but their Cartesian product may still be enormous.
The seven cache states are a protocol abstraction, not an enumeration of the implementation's bytes.
Chapter~\ref{ch:model-checking-algorithms} states the simulation obligation that connects such an abstraction to a concrete system.

The initial set is part of the specification.
Checking only $c_0$ proves a cold-start property; adding $c_4$ permits a snapshot taken during invalidation.
An unconstrained initial pin or metadata bit can introduce counterexamples before the first modeled instruction.
Conversely, an initial invariant that already assumes the property being checked can hide precisely the initialization bug of interest.

\section{Termination, errors, and total transitions}
\label{sec:ch45-deadlocks}

With a total relation, a \term{path} is an infinite sequence $\pi=s_0s_1\ldots$ satisfying $(s_i,s_{i+1})\in R$ for every $i\geq0$.
An initial path starts in $I$.
A finite execution that halts is represented by repeating its final state forever.
This convention lets one temporal semantics describe both termination and continued execution.
It does not assert that the runtime continues executing instructions after return.

We distinguish three terminal outcomes before adding loops: normal halt, explicit error, and unexpected stuckness.
For a raw successorless state, Algorithm~\ref{alg:ch45-totalize} records the outcome and adds the loop.
Existing labels are retained.
The input contract requires normal and error terminals to have no outgoing edges; a state with a real enabled transition is not a terminal merely because its name contains \code{done}.
For the cache, $c_5$ is a normal terminal and $c_6$ is an error terminal.

\begin{algorithm}[tbp]
\caption{Totalize a finite model while retaining terminal outcomes.}
\label{alg:ch45-totalize}
\begin{pseudocode}
\Require raw finite model $(S,I,R_0,AP,L)$, adjacency lists, terminal classification
\Ensure total $R$ and labels distinguishing halt, error, and deadlock
\State $R\gets R_0$
\ForAll{$s\in S$}
  \If{$\Succ_{R_0}(s)=\emptyset$}
    \If{$s$ is classified as normal termination}
      \State $L(s)\gets L(s)\cup\{\mathit{halt}\}$
    \ElsIf{$s$ is classified as explicit failure}
      \State $L(s)\gets L(s)\cup\{\mathit{error}\}$
    \Else
      \State $L(s)\gets L(s)\cup\{\mathit{deadlock},\mathit{error}\}$
    \EndIf
    \State $R\gets R\cup\{(s,s)\}$
  \EndIf
\EndFor
\State \Return $(S,I,R,AP\cup\{\mathit{halt},\mathit{error},\mathit{deadlock}\},L)$
\end{pseudocode}
\end{algorithm}

Scanning an adjacency-list representation costs $O(|S|+|R_0|)$ time and adds at most $|S|$ edges.
If successor counts are already available, the additional scan costs $O(|S|)$.
Label insertion is constant time for these three reserved Boolean flags.
The transformation does not change which old states are reachable, but it does specify the meaning of temporal statements after termination.
For example, $\mathbf{X}\mathit{halt}$ is true at a normal terminal and false at an error terminal lacking $\mathit{halt}$.
The property $\mathbf{F}\mathit{halt}$ rejects both a genuine infinite execution without halt and an error terminal that repeats forever.

Without totalization, a universal next operator at a dead end might be defined vacuously true, while an existential next operator is false.
Different finite-trace logics make other choices.
Neither interpretation should enter an implementation accidentally through an empty successor loop.
All temporal formulas and algorithms in these two chapters use the totalized, infinite-path convention unless a fairness restriction is explicitly stated.

\section{Safety, progress, and the meaning of a counterexample}
\label{sec:ch45-safety-liveness}

The cache invariant $\mathit{pinned}\Rightarrow\mathit{resident}$ holds in every baseline reachable state.
Adding $c_4\to c_6$ gives the finite witness
\[
 c_0,c_1,c_2,c_3,c_4,c_6.
\]
At its final state, the implication is false; no continuation can undo the fact that it was false there.
This is the characteristic shape of a \term{safety property}: each violating infinite trace has a finite prefix whose every infinite extension violates the property.
The definition concerns the property over labeled traces, not merely whether a particular checker prints a finite diagnostic.

A \term{liveness property} permits every finite trace prefix to be extended to a satisfying infinite trace.
For example, a request must eventually receive a response.
An arbitrary finite delay alone does not establish a violation, because a response can still occur later.
The cache requirement that entry into $c_4$ eventually leads to $c_5$ fails on
\[
 c_0c_1c_2c_3(c_4)^\omega,
\]
where $\omega$ means infinite repetition.
There is no finite deadline in the requirement.
An implementation that reports failure merely after a large number of iterations has checked a bounded-response property instead.

Safety and liveness are semantic classifications, not a list of keywords.
Properties can combine them.
A requirement that reclamation is safe and eventually happens has both an invariant part and a progress part.
Moreover, the model can make a liveness failure inevitable after a finite prefix, for example by entering an absorbing error state.
That observation does not turn the unrestricted trace property \enquote{eventually halt} into a safety property.
It uses information about which extensions the particular model allows.

Not all safety properties are state invariants over the original state variables.
Suppose a deoptimization protocol emits mutually exclusive events $\mathit{fail}$, $\mathit{rebuild}$, and $\mathit{commit}$, with $\mathit{none}$ for steps without an event.
We want to forbid a commit after a guard failure until a reconstruction event has occurred.
A three-state monitor remembers the relevant history: $C$ means clear, $P$ means pending reconstruction, and $B$ means violation.
From either nonbad state, a failure enters $P$ and a reconstruction enters $C$; a commit from $P$ enters $B$; all other events preserve the state.
State $B$ is absorbing.
The event sequence $\mathit{fail},\mathit{none},\mathit{commit}$ yields monitor states $P,P,B$.
The sequence $\mathit{fail},\mathit{rebuild},\mathit{commit}$ yields $P,C,C$.

Algorithm~\ref{alg:ch45-monitor} constructs the product with this deterministic monitor.
The initial states are just before the first event, and each transition carries one event $e(s,t)$.
For ordinary state-labeled Kripke input, an event can be stored in the target state; if two incoming transitions have different events, split that state or retain an explicit event component.
The product is again a state model, with $\mathit{violation}$ labeling states whose monitor component is $B$.

\begin{algorithm}[tbp]
\caption{Turn a finite-history deoptimization obligation into reachability.}
\label{alg:ch45-monitor}
\begin{pseudocode}
\Require total finite model $(S,I,R)$, edge events $e$, monitor $Q=\{C,P,B\}$ and $\delta$
\Ensure reachable product states and edges, initialized before the first event
\State $D\gets\{(s,C)\mid s\in I\}$; $W\gets$ queue containing $D$; $R_P\gets\emptyset$
\While{$W\neq\emptyset$}
  \State remove $(s,q)$ from $W$
  \ForAll{$t\in\Succ(s)$}
    \State $q'\gets\delta(q,e(s,t))$; $R_P\gets R_P\cup\{((s,q),(t,q'))\}$
    \If{$(t,q')\notin D$}
      \State add $(t,q')$ to $D$ and $W$
    \EndIf
  \EndFor
\EndWhile
\State \Return $D,R_P$, with bad states $\{(s,B)\in D\}$
\end{pseudocode}
\end{algorithm}

For a monitor with $|Q|$ states, the product has at most $|S||Q|$ states and $|R||Q|$ edges.
With constant-time indexed membership and monitor transitions, construction takes $O(|Q|(|S|+|R|))$ time.
The bound concerns this finite deterministic monitor, not translation of an arbitrary temporal formula.
The invariant that the monitor component equals the fold of $\delta$ over the path's events proves that reaching $B$ is equivalent to violating the stated ordering rule.
For multiple concurrent deoptimizations, one global pending bit is insufficient: a reconstruction for one activation must not clear another activation's obligation.
The identity of the obligation belongs in the model.

\section{Linear-time formulas quantify over complete paths}
\label{sec:ch45-ltl}

\term{Linear temporal logic} (LTL) describes one infinite path at a time.
We write $\pi,i\models p$ when $p\in L(s_i)$.
Boolean connectives have their ordinary meanings at position $i$.
The temporal operators are defined by
\begin{align*}
 \pi,i\models\mathbf{X}\varphi
 &\quad\Longleftrightarrow\quad \pi,i+1\models\varphi,\\
 \pi,i\models\varphi\,\mathbf{U}\,\psi
 &\quad\Longleftrightarrow\quad
 \exists j\geq i:\ \pi,j\models\psi
 \text{ and }\forall k\in[i,j):\pi,k\models\varphi.
\end{align*}
The until operator is \emph{strong}: the $\psi$ position must exist.
The interval before it may be empty, so $\psi$ now is sufficient.
Define $\mathbf{F}\varphi=\mathit{true}\,\mathbf{U}\,\varphi$ and $\mathbf{G}\varphi=\neg\mathbf{F}\neg\varphi$.
Thus \enquote{eventually} includes the current position, and \enquote{always} includes it too.

The model satisfies an LTL formula when \emph{every} initial path satisfies it:
\[
 K\models_{\mathrm{LTL}}\varphi
 \quad\Longleftrightarrow\quad
 \forall\pi\text{ starting in }I:\ \pi,0\models\varphi.
\]
Universal path quantification is the model-checking convention, not an operator that LTL permits us to insert inside a path formula.
A single satisfying execution does not establish an LTL property of the model.

For the baseline cache,
\[
 \mathbf{G}(\mathit{published}\Rightarrow\mathit{metadata})
 \quad\text{and}\quad
 \mathbf{G}(\mathit{pinned}\Rightarrow\mathit{resident})
\]
both hold by inspection of all reachable labels.
The response property $\mathbf{G}(\mathit{draining}\Rightarrow\mathbf{F}\mathit{freed})$, where $\mathit{draining}$ labels only $c_4$, fails on the infinite waiting path above.
By contrast, the stronger bounded rule $\mathbf{G}(\mathit{draining}\Rightarrow\mathbf{X}\mathit{freed})$ already fails on a single waiting step at $c_4$.
The formulas ask different engineering questions.

Atomicity affects formulas containing $\mathbf{X}$.
Splitting publication into two steps changes what \enquote{next} means, even if it introduces no new observation beyond repeating a label.
LTL without $\mathbf{X}$ is invariant under stuttering of labels, but this fact alone does not justify collapsing a concurrent protocol: the collapse must still preserve the relevant paths~\cite{baier-2008-model-checking}.
In particular, hiding an intermediate state where published code lacks metadata hides the actual invariant violation rather than a harmless stutter.

\section{Branching-time formulas retain choices at each state}
\label{sec:ch45-ctl}

A pass manager may always retain a route to a consistent analysis cache without guaranteeing that it ever takes that route.
Figure~\ref{fig:ch45-pass} models this distinction.
State $a$ begins a transformation, $b$ has invalidated analysis information and may retry, $c$ has rebuilt valid information, and $d$ has finished.
The propositions $\mathit{valid}$ and $\mathit{done}$ label $\{c,d\}$ and $\{d\}$, respectively.
The only initial state is $a$.

\begin{figure}[tbp]
\centering
\begin{tikzpicture}[node distance=15mm]
\node[basic block] (a) {$a$: transform};
\node[basic block,right=of a] (b) {$b$: invalid};
\node[basic block,right=of b] (c) {$c$: valid};
\node[basic block,right=of c] (d) {$d$: done};
\draw[cfg edge] ([xshift=-8mm]a.west) -- (a);
\draw[cfg edge] (a) -- (b);
\draw[cfg edge] (b) -- node[above,font=\scriptsize] {rebuild} (c);
\draw[cfg edge] (c) -- (d);
\draw[cfg edge] (b) to[loop above] node[above,font=\scriptsize] {retry} (b);
\draw[cfg edge] (d) to[loop above] (d);
\end{tikzpicture}

\caption{A rebuild remains available at every retry, but the retry loop can postpone it forever.}
\label{fig:ch45-pass}
\end{figure}

\term{Computation tree logic} (CTL) evaluates state formulas.
Each temporal operator is paired with an existential path quantifier $\mathrm{E}$ or a universal path quantifier $\mathrm{A}$.
For example, $s\models\mathrm{EX}\varphi$ means that some successor satisfies $\varphi$, while $s\models\mathrm{AX}\varphi$ means that every successor does.
The formula $\mathrm{E}[\varphi\,\mathrm{U}\,\psi]$ requires one path with a future $\psi$ state and $\varphi$ at all preceding positions.
The formula $\mathrm{EG}\varphi$ requires one path whose every state satisfies $\varphi$.
The universal variants require the corresponding condition on every path from the state.
The model satisfies a CTL formula when every state in $I$ satisfies it.

Crucially, the operands of these operators are themselves state formulas, so path quantifiers can nest.
The formula $\mathrm{AG}\,\mathrm{EF}\mathit{done}$ says that from every reachable state there exists a path to completion.
It is true at $a$: from $b$ choose $c$, then $d$; from $c$ choose $d$; and at $d$ the eventual target is already true.
But $\mathrm{AF}\mathit{done}$ is false at $a$, because $ab^\omega$ never completes.
The existential choice inside $\mathrm{AG}\,\mathrm{EF}$ can be different at each visited state.
It does not oblige the outer path to follow that choice.

Table~\ref{tab:ch45-ctl-truth} records the complete calculation.
For $\mathrm{EG}\neg\mathit{done}$, states $a$ and $b$ can reach and remain on the retry loop without seeing $d$; state $c$ must reach $d$ on its next step.
For $\mathrm{AX}\mathit{valid}$, $b$ fails because one of its successors is $b$ itself.
These are state-by-state truth values, not frequencies or probabilities.

\begin{table}[tbp]
\centering
\caption{A path to completion exists everywhere, although completion is not inevitable everywhere.}
\label{tab:ch45-ctl-truth}
\begin{tabular}{lccccc}
\toprule
State & $\mathrm{EF}\mathit{done}$ & $\mathrm{AF}\mathit{done}$ & $\mathrm{AG}\,\mathrm{EF}\mathit{done}$ & $\mathrm{EG}\neg\mathit{done}$ & $\mathrm{AX}\mathit{valid}$ \\
\midrule
$a$ & T & F & T & T & F \\
$b$ & T & F & T & T & F \\
$c$ & T & T & T & F & T \\
$d$ & T & T & T & F & T \\
\bottomrule
\end{tabular}
\end{table}

LTL $\mathbf{F}\mathit{done}$ under universal model semantics agrees here with CTL $\mathrm{AF}\mathit{done}$.
That agreement for one formula is not an equivalence of the logics.
CTL can express the continuing availability of alternatives, as in $\mathrm{AG}\,\mathrm{EF}\mathit{done}$, which universal LTL cannot generally express.
Conversely, LTL can express path properties such as eventual permanent stability, $\mathbf{F}\mathbf{G}p$, that have no equivalent CTL formula in general~\cite{huth-2004-logic}.
Choose the logic to match the intended requirement rather than translating words by matching operator names.

\section{Fairness is an assumption about allowed paths}
\label{sec:ch45-fairness}

The retry loop is a real behavior of the transition relation.
We may exclude it only by stating a justified environment or scheduler assumption.
A \term{fairness constraint} selects a subset of infinite paths; it does not delete an arbitrary edge or declare every infinite execution acceptable.
For sets $J_1,\ldots,J_m\subseteq S$, a generalized B\"uchi fairness condition requires each $J_i$ to be visited infinitely often.
Write $\mathcal{F}(\pi)$ for that condition.
Fair LTL checking asks whether every initial path satisfying $\mathcal{F}$ satisfies the property.

On the pass model, choose $J_1=\{c,d\}$.
The path $ab^\omega$ is unfair, whereas $ab^nc d^\omega$ is fair for every finite $n\geq1$.
Under this restriction, $\mathbf{F}\mathit{done}$ holds.
The model has not acquired a finite upper bound on retries: the fair paths still permit arbitrarily large $n$.
Because $c$ is followed by $d$ and $d$ repeats, this particular fairness condition says that the rebuild is eventually taken.
It must be justified by the scheduler or by a separate progress argument; it cannot serve as an independent proof of the same progress obligation.

For actions, \term{weak fairness} requires an action that remains continuously enabled from some point onward to occur infinitely often.
Equivalently, it cannot remain enabled forever while being ignored forever.
\term{Strong fairness} imposes the corresponding obligation when the action is enabled infinitely often, even if enabling is intermittent.
If $\mathit{enabled}_a$ and $\mathit{taken}_a$ are encoded as observations with a consistent step convention, the respective conditions can be written
\[
 (\mathbf{F}\mathbf{G}\mathit{enabled}_a\Rightarrow\mathbf{G}\mathbf{F}\mathit{taken}_a),
 \qquad
 (\mathbf{G}\mathbf{F}\mathit{enabled}_a\Rightarrow\mathbf{G}\mathbf{F}\mathit{taken}_a).
\]
For a one-shot action, taking it may disable it permanently, making the antecedent false on the completed path.
In the cache model the release action is continuously enabled while waiting at $c_4$.
Weak fairness for release excludes $c_4^\omega$; it says nothing about a runtime in which the user can permanently block before release becomes enabled.

Fair CTL also restricts its path quantifiers, including nested ones.
An existential witness must have a fair infinite continuation, even when its until target is reached immediately.
If a state has no fair continuation, universal fair-path formulas can become vacuously true there.
A useful accompanying check is that every intended initial state has at least one fair continuation.
The ordinary CTL algorithms in Chapter~\ref{ch:model-checking-algorithms} do not silently apply these restrictions.

For compiler engineering, record the state variables, step boundaries, initial condition, terminal classification, and fairness assumptions beside the property.
A successful check of the seven-state cache protocol does not establish a memory-ordering theorem for a lock-free implementation.
A successful fair progress check does not establish that the operating system schedules a blocked worker.
These are boundaries of the mathematical claim, and they identify the additional models or proofs needed before using the result to change an optimizer or runtime.

\section{Exercises}
\label{sec:ch45-exercises}

\begin{enumerate}
\item\label{ex:ch45-cache-witness}
Add only the dashed edge $c_4\to c_6$ to the cache model.
Find a shortest violation of $\mathbf{G}(\mathit{pinned}\Rightarrow\mathit{resident})$, count its transitions, and decide whether the publication-metadata invariant still holds.

\item\label{ex:ch45-terminal-policy}
A raw model has initial state $u$, edges $u\to h$ and $u\to z$, and no other edges.
State $h$ is a normal terminal; $z$ is unclassified stuckness.
Apply Algorithm~\ref{alg:ch45-totalize} and evaluate $\mathrm{EF}\mathit{halt}$, $\mathrm{AF}\mathit{halt}$, and $\mathrm{AG}\neg\mathit{error}$ at $u$.

\item\label{ex:ch45-monitor-history}
Run the deoptimization monitor on $\mathit{fail},\mathit{rebuild},\mathit{fail},\mathit{commit}$.
Explain why a Boolean flag meaning \enquote{some reconstruction has occurred} is insufficient.
Does adding the event $\mathit{rebuild}$ after the final commit repair this trace?

\item\label{ex:ch45-branching}
In the pass model add a new state $e$ with only $e\to e$, add $b\to e$, and label $e$ with neither $\mathit{done}$ nor $\mathit{valid}$.
Compute $\mathrm{EF}\mathit{done}$ and $\mathrm{AG}\,\mathrm{EF}\mathit{done}$ at all five states without fairness.
Explain the different effects of an optional failed route and a permanently lost route to completion.

\item\label{ex:ch45-fairness-gap}
An infinite execution alternates states $u,v,u,v,\ldots$.
An action $a$ is enabled exactly at $u$ and is never taken.
Does the execution satisfy weak fairness for $a$? Strong fairness?
Explain why the answer changes if $a$ is enabled at both states.
\end{enumerate}

\section{Further reading}
\label{sec:ch45-reading}

Baier and Katoen develop transition-system semantics, safety and liveness, and fairness with careful separation of the different assumptions~\cite{baier-2008-model-checking}.
Huth and Ryan give a complementary treatment of LTL and CTL that is useful when checking the scope of nested path quantifiers~\cite{huth-2004-logic}.
The \emph{Handbook of Model Checking} connects these semantic choices to automata, symbolic representations, and verification tools~\cite{clarke-2018-handbook}.
